% 金星（综述）
% license CCBYSA3
% type Wiki

本文根据 CC-BY-SA 协议转载翻译自维基百科\href{https://en.wikipedia.org/wiki/Venus}{相关文章}。

金星是距离太阳第二近的行星。由于其轨道最接近地球，且两者同为类地岩石行星，并在大小、质量和表面引力上几乎完全一致，因此金星常被称为地球的“孪生”或“姐妹”行星。然而，金星在许多方面与地球显著不同，尤其是它没有液态水，其大气层远比太阳系中其他任何岩石天体都更加厚重和致密。金星大气主要由二氧化碳组成，并覆盖着一层浓密的硫酸云层，环绕整个行星。在平均地表高度处，大气温度达到 737\,K（464\,°C；867\,°F），压力为地球海平面的 92 倍，使其最低层大气呈现超临界流体状态。从地球观测时，金星呈现为类似恒星的亮点，比天空中任何其他天然光源都更明亮[22][23]，并且作为一颗内行星，它总是出现在接近太阳的方向，要么作为最明亮的“启明星”，要么作为“长庚星”。

金星与地球的轨道使两者在约 1.6 年的会合周期中彼此接近。在此过程中，金星会比其他任何行星更接近地球；相较之下，由于水星的轨道更靠近太阳，因此它在自身轨道路径中整体上比其他行星更接近地球。在从地球出发的行星际空间飞行中，金星常被用作引力助推的中转点，从而提供更快速且更经济的航程。金星没有卫星，并以非常缓慢的逆行方式绕其自转轴旋转，这被认为是太阳潮汐锁定效应与金星巨大大气层的受热差异相互竞争的结果。因此，金星的一昼夜长达 116.75 个地球日，而其一个太阳年为 224.7 个地球日。

金星具有微弱的磁层；由于缺乏内部发电机，其磁场是由太阳风与大气相互作用所诱导产生的。在内部结构上，金星具有核心、地幔和地壳。内部热量通过活跃的火山活动逸出[24][25]，从而导致地表更新，而不是板块构造。金星在其早期历史中可能曾拥有液态地表水及宜居环境[26][27]，但由于失控的温室效应，所有水分最终蒸发，使金星演变为当前的状态[28][29][30]。在云层高度处的大气条件是整个太阳系中与地球最为相似的，这些条件被认为可能有利于金星生命的存在，并在 2020 年发现了潜在的生物标志物，从而推动了新的研究与探测任务。

在人类历史上，全世界的人们都曾观测过金星，并使其在众多文化中具有特殊的重要性。随着望远镜的出现，金星的相位变得可辨识，并在 1613 年被呈现为推翻当时占主导地位的地心模型、支持日心模型的决定性证据。1961 年，“金星 1 号”首次造访金星，飞掠行星并实现了人类第一次行星际飞行。1962 年，第二次行星际任务“水手 2 号”带回了首批来自金星的数据。1967 年，首个行星际撞击器“金星 4 号”抵达金星，随后在 1970 年由着陆器“金星 7 号”成功登陆。截至 2025 年，“太阳轨道飞行器”正前往于 2026 年飞掠金星，而下一项计划发射前往金星的任务是“金星生命探测器”，预定同样于 2026 年发射。







